% Source fingerprint: f845e21e27a6874ee002bb345436ca13545a1e22a1fa5ebbe9c9e09117f3263b
% T17: raw TeX cells preserved; no numeric highlighting.
\begin{table}[htbp]
\centering\small\renewcommand{\arraystretch}{1.18}\setlength{\tabcolsep}{4pt}
\caption[Fredholm 可解性的检查项目]{Fredholm 可解性的检查项目。核有限维、值域闭和指标零是三个不同结论。一般 Banach 空间先用对偶配对，不能无条件借用正交语言。}
\label{b1:tab:visual-t17}
\begin{tabularx}{\linewidth}{@{}>{\raggedright\arraybackslash}p{3.1cm}>{\raggedright\arraybackslash}p{4.7cm}>{\raggedright\arraybackslash}X@{}}
\toprule
算子情形 & 已获得的结构 & 方程的相容条件 \\
\midrule
一般有界线性算子 & 核闭；值域不一定闭 & 不能因存在近似解就断言右端属于值域 \\
闭值域算子 & \(X/\ker A\) 上有由 \(\|Ax\|\) 控制的估计 & 闭值域定理把值域识别为相应对偶核的零化子 \\
\(I-K\)，\(K\) 紧 & 有限维核、闭值域、有限维余核，指标为零 & 单射等价于满射；不要求 \(K\) 自伴或可以对角化 \\
Hilbert 空间中的闭值域 & 可用伴随和正交补表示 & \(Ax=y\) 可解当且仅当 \(y\perp\ker A^*\) \\
紧自伴谱方法 & 沿正交特征向量逐分量处理 & 零特征值处有相容条件；小非零特征值还要求解系数可平方求和 \\
\bottomrule
\end{tabularx}
\end{table}
